\documentclass[10pt,a4paper]{amsart}
\usepackage[margin=19mm]{geometry}
\usepackage{amsmath,amssymb,mathtools}
\usepackage[scr=rsfs,cal=euler]{mathalfa}
\usepackage{xfrac}
\usepackage[unicode,hidelinks,bookmarks=false]{hyperref}
\newcommand{\ie}{i.e.\ }
\newcommand{\cf}{cf.\ }
\newcommand{\resp}{respectively\ }
\newcommand{\Subsection}{\subsection}
\providecommand{\nobreakdash}{\nobreak\hbox{-}}
\setlength{\emergencystretch}{2em}
\multlinegap0pt
\usepackage[shortlabels]{enumitem}
\usepackage{amssymb}
\usepackage{mathtools}
\usepackage{xcolor}
\usepackage{tikz}
\usepackage{float}
\newtheorem{lemma}{Lemma}
\newtheorem{corollary}{Corollary}
\newtheorem{theorem}{Theorem}
\newcommand{\dist}{\operatorname{dist}}
\title{Exact Leaf Powers on Cycles, Ladders, Crowns, and Multipartite Block Graphs}
\author{Peng Li, Yangjing Long}
\date{Source: arXiv:2606.02271v1 (CC BY 4.0)}
\begin{document}
\maketitle

\noindent\emph{Source and licence.} Exact Leaf Powers on Cycles, Ladders, Crowns, and Multipartite Block Graphs. Version arXiv:2606.02271v1 (CC BY 4.0), distributed by arXiv under the Creative Commons Attribution 4.0 International licence, http://creativecommons.org/licenses/by/4.0/, as recorded in the arXiv licence metadata for this exact version and shown on the abstract page https://arxiv.org/abs/2606.02271v1. Full licence notice: \texttt{licenses/arxiv-2606.02271-CC-BY-4.0.txt}.\par\smallskip

\noindent\textit{Editorial scope.} Counted unit: the article's theorem thm:ladder (source0.tex lines 362--367). The article's complete original proof of it is delivered with it (source0.tex lines 369--414, one \texttt{proof} environment of 46 lines). No mathematics is written, repaired or re-derived: the statement and the proof are the article's own lines, in the article's own notation, and no retained line is substituted or re-flowed.  Retained uncounted as the source-local context the proof needs: lines 230--232; lines 238--240; lines 337--351.  Nothing was omitted from any retained span, so the package carries no omission record.  No retained line contains a citation, so the wrapper carries no bibliography and no bibliography entry.  No retained line contains a figure environment, an image-inclusion command or drawing code, so no image is delivered and the package declares no asset dependency.  Editorial formatting only, in addition: an AMS \texttt{amsart} preamble that restates the article's own declarations and macros used by the retained lines, namely the environments \texttt{lemma}, \texttt{corollary}, \texttt{theorem} on separate counters, together with the standard packages those lines need.  The retained environments print in this document as Lemma 1, Corollary 1, Theorem 1 in order of appearance, whereas the article prints the same items with the numbering of its own full document.  The complete native source of the article as distributed in the arXiv e-print for this exact version is bundled unchanged as \texttt{sources/201890/source0.tex}.
\begin{lemma}[support reduction]\label{lem:support-reduction}
A graph $G$ is an exact $5$-leaf power if and only if it has a support-tree distance-three representation.
\end{lemma}

\begin{corollary}[induced-subgraph closure]\label{cor:induced-closure}
For $k\ge2$, every induced subgraph of an exact $k$-leaf power is an exact $k$-leaf power.
\end{corollary}

\begin{lemma}[domino support form]\label{lem:domino-support-form}
Let $D$ be the domino with bipartition
\[
 A=\{a_0,a_1,a_2\},\qquad B=\{b_0,b_1,b_2\},
\]
and edges $a_i b_j$ exactly when $|i-j|\le1$.  In every support-tree distance-three representation of $D$, exactly one of the following two symmetric alternatives holds, up to reversing the order $0,1,2$.
\begin{enumerate}[label=\textup{(\alph*)}]
\item There is a vertex $c$ adjacent to $s_{a_0},s_{a_1},s_{a_2}$, a vertex $d$ adjacent to $c$ and to $s_{b_1}$, and the two boundary supports satisfy
\[
 s_{b_0}s_{a_2}\in E(R),\qquad s_{b_2}s_{a_0}\in E(R).
\]
We call this the $A$-oriented support form.
\item The symmetric $B$-oriented form holds, with the roles of $A$ and $B$ interchanged.
\end{enumerate}
\end{lemma}

\begin{theorem}[ladders]\label{thm:ladder}
For $t\ge1$, the ladder $L_t$ is an exact $5$-leaf power if and only if
\[
 t\le2.
\]
\end{theorem}

\begin{proof}
By Lemma~\ref{lem:support-reduction}, it is enough for the positive cases to give support-tree distance-three representations.

For $L_1=K_{2,2}$, take a path $a b$ and attach the two $X$-support vertices $s_{x_0},s_{x_1}$ to $a$ and the two $Y$-support vertices $s_{y_0},s_{y_1}$ to $b$.  Then every $s_{x_i}$ is at distance three from every $s_{y_j}$, while the same-side distances are two.  Hence the exact distance-three graph is $K_{2,2}=L_1$.

For $L_2$, take a support tree with vertices
\[
 s_{x_0},s_{x_1},s_{x_2},\quad s_{y_0},s_{y_1},s_{y_2},\quad c,d
\]
and edges
\[
 cs_{x_0},\ cs_{x_1},\ cs_{x_2},\ cd,\ ds_{y_1},\ s_{x_2}s_{y_0},\ s_{x_0}s_{y_2}.
\]
Then
\[
 \dist(s_{x_i},s_{y_j})=3
 \quad\Longleftrightarrow\quad |i-j|\le1,
\]
for $i,j\in\{0,1,2\}$, and no same-side pair is at distance three.  Therefore the represented graph is $L_2$.

It remains to exclude $t\ge3$.  By Corollary~\ref{cor:induced-closure}, it is enough to exclude $L_3$.  Suppose, for a contradiction, that $L_3$ has a support-tree distance-three representation $R$.

Let $D_L$ be the induced domino on
\[
 \{x_0,x_1,x_2,y_0,y_1,y_2\},
\]
and let $D_R$ be the induced domino on
\[
 \{x_1,x_2,x_3,y_1,y_2,y_3\}.
\]
The restriction of $R$ to the supports of $D_L$ is a support representation of a domino, and similarly for $D_R$.  Hence Lemma~\ref{lem:domino-support-form} applies to both.

Assume first that $D_L$ is $X$-oriented.  Thus $s_{x_0},s_{x_1},s_{x_2}$ have a common neighbour $c$, while the boundary support $s_{y_0}$ is adjacent to $s_{x_2}$ and the boundary support $s_{y_2}$ is adjacent to $s_{x_0}$, up to reversing the left domino.  The two shared $Y$-supports $s_{y_1}$ and $s_{y_2}$ are then at distance four in the left domino support form.  Hence $D_R$ cannot be $Y$-oriented: in a $Y$-oriented support form the shared supports $s_{y_1}$ and $s_{y_2}$ would have a common neighbour and hence distance two, contradicting uniqueness of the path between them in the tree.  Thus $D_R$ is also $X$-oriented.

In the $X$-oriented form of $D_R$, the supports $s_{x_1}$ and $s_{x_2}$ have a common neighbour.  In a tree, two vertices at distance two have a unique common neighbour; since their common neighbour in $D_L$ is $c$, the common neighbour in $D_R$ is again $c$.  Therefore $s_{x_3}$ is adjacent to $c$.  But $s_{y_0}$ is adjacent to $s_{x_2}$ and $s_{x_2}$ is adjacent to $c$, so
\[
 \dist_R(s_{x_3},s_{y_0})=3.
\]
This would make $x_3y_0$ an edge of the represented graph, whereas $|3-0|>1$ and hence $x_3y_0\notin E(L_3)$.  This contradiction excludes the case that $D_L$ is $X$-oriented.

The case in which $D_L$ is $Y$-oriented is symmetric.  Then $D_R$ is forced to be $Y$-oriented, the shared neighbour of $s_{y_1}$ and $s_{y_2}$ is the same in both dominoes, and the boundary path forces
\[
 \dist_R(s_{x_0},s_{y_3})=3,
\]
contradicting the nonedge $x_0y_3$.  Thus $L_3$ is not an exact $5$-leaf power.  Since $L_t$ contains an induced $L_3$ for every $t\ge3$, no such $L_t$ is an exact $5$-leaf power.
\end{proof}

\end{document}
